\documentclass[11pt]{article}
\usepackage[T1]{fontenc}
\usepackage[margin=25mm]{geometry}
\usepackage{graphicx,booktabs,tabularx,placeins,amsmath,hyperref,microtype}
\hypersetup{hidelinks}
\newcommand{\ArchiveCount}{20}
\newcommand{\MeasuredArchives}{20}
\newcommand{\MeasuredPlatforms}{6}
\newcommand{\MaxNullSlice}{28.1}
\newcommand{\MaxUnalignedNull}{60.4}
\newcommand{\PublicRuns}{22}
\newcommand{\FixtureCount}{4}
\newcommand{\PhaseCells}{10080}
\newcommand{\SimCohorts}{1935360}
\newcommand{\RepeatCount}{96}
\newcommand{\MemoryCells}{1000000}
\newcommand{\MemoryFeatures}{64}
\newcommand{\MemoryPatients}{10}
\newcommand{\RSSGiB}{0.278}
\newcommand{\DensityPct}{2.96}
\newcommand{\ADUnits}{6}
\newcommand{\ADFeatures}{32}
\newcommand{\ADFloor}{0.10}
\newcommand{\ADSpot}{87.5}
\newcommand{\ADPatient}{0.0}
\newcommand{\SOTARuns}{77}
\newcommand{\SOTAMethods}{6}

\title{Spatial Audit: a generalizable five-dimensional statistical and experimental-design audit framework for multi-patient spatial omics}
\author{Chaohao Liu\thanks{Independent Researcher. Corresponding author email: \textit{[insert email before submission]}. ORCID: 0009-0003-4341-2494. Code and interactive observatory: \url{https://github.com/XG221B/spatial-audit} and \url{https://xg221b.github.io/spatial-audit/}.}}
\date{}
\newcommand{\auditfigure}[3]{\begin{figure}[htbp]\centering\includegraphics[width=.83\linewidth]{figures/#1.pdf}\caption{\textbf{#2} #3}\label{fig:#1}\end{figure}}
\begin{document}
\maketitle

\begin{abstract}
\noindent\textbf{Objective:} In clinical spatial omics, a tissue section can contribute thousands of measurements while adding only one donor to a comparison. We developed Spatial Audit to help investigators decide whether a proposed biomarker contrast supports donor-level testing, warrants descriptive reporting, or requires additional sampling.

\noindent\textbf{Materials and Methods:} We audited 20 primary human archives across six platforms. For each contrast, we assessed the resolution of donor-label permutations, alignment between clinical group and acquisition conditions, and sensitivity to pooling observations across donors. Where tissue coverage and metadata permitted, we also measured the fraction of spatial variation shared with anatomical or instrument covariates at specified wavelengths and examined sensitivity to short-range noise. These checks used calibrated coordinates and donor-linked tissue sections. We compared native statistical methods on registered panels and randomized labels within pretreatment breast cancer response groups.

\noindent\textbf{Results:} All 20 archives supported permutation resolution and pooled versus donor comparisons. None supplied shared instrument stage coordinates. Eleven supported a covariate projection at one or more wavelengths, and nine supported noise assessment for at least one feature. Conditional on the observed donors and 32 features, donor testing after Benjamini--Hochberg correction called 43/1,600 feature-splits in residual disease (4/50 splits with calls) and 27/1,600 in pathological complete response (1/50). Pooled cells yielded 1,493/1,600 and 1,521/1,600, respectively (50/50 splits in each group). In cortex, three donors per arm imposed a two-sided permutation floor of 0.10 at both panel sizes.

\noindent\textbf{Conclusion:} Reporting which contrasts the available donors, slides and tissue footprints can support helps clinical research teams prioritize findings for independent validation. An unresolved acquisition effect or an unsupported spatial scale can change what is defensible even when a pooled analysis identifies many candidate markers.

\end{abstract}
\noindent\textbf{Keywords:} Spatial transcriptomics; Mass spectrometry imaging; Biomedical informatics; Pseudoreplication; Confounding audit; Reproducibility; Permutation testing.

\begin{table}[htbp]\centering
\caption{Statement of Significance}\label{tab:significance}
\begin{tabularx}{\linewidth}{@{}p{.23\linewidth}X@{}}\toprule
\textbf{Problem or issue} & A biomarker contrast can appear precise when many cells come from a few donors or when clinical groups occupy different slides.\\
\textbf{What is already known} & Local quality checks identify unreliable measurements. Donor-aware tests address biological replication, and spatial variability tests identify patterns within tissue.\\
\textbf{What this paper adds} & Across 20 public archives, the audit shows which donor contrasts and spatial scales can be assessed. The results inform decisions about descriptive reporting, donor-level testing and additional sampling.\\
\bottomrule\end{tabularx}
\end{table}

\FloatBarrier
\section{Introduction}
Spatially resolved molecular profiling---spanning capture-array transcriptomics (Visium, Stereo-seq), targeted in situ imaging (Xenium, CosMx SMI, MERFISH), and matrix-assisted laser desorption/ionization mass spectrometry imaging (MALDI-MSI)---has rapidly entered clinical and translational research. Across oncology, neurodegeneration, immunology, and cardiopulmonary disease, multi-patient spatial cohorts are increasingly used to identify disease-associated markers, treatment-response programs, and tissue microenvironments. Yet translating high-dimensional spatial assays into reproducible biomedical findings confronts a fundamental statistical asymmetry: a single study may collect hundreds of thousands of spots, cells, or pixels, while the number of independent patients per clinical arm often remains between three and fifteen.

When downstream pipelines treat spatially correlated within-patient observations as independent replicates, test statistics become severely anti-conservative and nominal $p$-values collapse toward zero even under the null hypothesis \cite{pseudorep}. Conversely, when inference is properly anchored at the biological individual, small donor counts impose a hard combinatorial floor on exact permutation $p$-values regardless of how densely each tissue section is sampled. Beyond sample size, clinical spatial comparisons face three structural confounders that routinely interact:
\begin{enumerate}
\item \textbf{Acquisition-design aliasing:} Clinical group labels can align with slide assignment, processing batch, or physical instrument stage coordinates, rendering biological differences partially or wholly unidentifiable from technical effects.
\item \textbf{Multi-scale composition and instrument confounding:} Apparent spatial patterning at a given physical wavelength ($\mu\mathrm{m}$) can be explained by macro-anatomical region occupancy, cell-type composition, or---in MALDI-MSI---total ion current (TIC) and raster drift rather than within-compartment molecular regulation.
\item \textbf{Abundance-linked short-range noise:} Pixel- and count-level high-frequency variance scales with baseline abundance, so shifts in mean expression between clinical arms can create apparent fine-scale spatial energy differences even in the absence of latent spatial structure.
\end{enumerate}

\subsection{Background and related work}
Spatial quality control assesses whether individual spots, segmented cells or tissue regions yield interpretable measurements. A section can pass these checks even when all cases were processed on one slide and all controls on another. Donor-aware differential expression addresses replication \cite{pseudorep,edger,deseq2}. Domain-interaction and within-section variability methods address the location of molecular patterns \cite{despace2,nnsvg,sparkx,cardinal}. Spatial Audit asks which clinical contrasts and physical scales the sampling design supports before those analyses are interpreted (Table~\ref{tab:significance}).

Consider a team selecting a cortical marker for follow-up. With three case and three control donors, the smallest two-sided donor permutation probability is 0.10. The team can report donor effect estimates descriptively. If each arm occupies its own slide, additional donors processed on shared slides would help distinguish a group effect from acquisition. Once replication and slide assignment support a donor test, the proposed wavelength still matters. A coarse pattern shared with cortical layers motivates a layer-specific comparison, while a wavelength below the array's resolution calls for finer sampling before interpreting local variation.

\section{Materials and Methods}
\subsection{Data contract and study selection}
\textit{Spatial Audit} operates on standardized cohorts supplied either via portable JSON manifests (referencing NPZ or backed HDF5/H5AD files) or directly from a single AnnData/H5AD file through the \texttt{quick-run} command and \texttt{audit\_anndata} Python API. Every observation row carries calibrated two-dimensional physical coordinates in micrometres ($\mu\mathrm{m}$), a molecular abundance vector, optional spatial covariates $C$, and section-level metadata linking each tissue section to a biological unit (\texttt{patient\_id}), clinical group, and optional \texttt{slide\_id}, \texttt{batch}, and physical instrument \texttt{stage} coordinates.

Public archives in our benchmark catalog were selected for (i) spatially resolved human measurements with recoverable physical scale, (ii) identifiable biological groups or a clearly scoped tissue contrast, and (iii) a directly retrievable official accession. Every measured public report preserves SHA-256 input receipts, coordinate calibration notes, donor crosswalks, and strict RFC-8259 JSON serialization (representing undefined quantities as \texttt{null}). Public outputs and deterministic synthetic fixtures are strictly segregated.

\subsection{Dimension 1: Exchangeability and permutation resolution}
Independent-group audits require a unique group label per biological individual and aggregate multiple sections within each donor before group inference. For independent group sizes $n_a$ and $n_b$, there are $M = \binom{n_a + n_b}{n_a}$ legal label allocations. We define two-sided permutation probabilities by doubling the smaller inclusive single-tail probability and capping at one, yielding an exhaustive theoretical floor of $p_{\min} = \min(1, 2/M)$. Across a registered panel of $m$ features at significance level $\alpha$ (default $0.05$, user-configurable), the smallest feasible strict Benjamini--Hochberg (BH) rejection rank is $\lfloor m p_{\min}/\alpha \rfloor + 1$ \cite{bh}. When $M$ exceeds the enumeration budget, reproducible Monte Carlo allocations are drawn with the observed assignment included in both numerator and denominator, and theoretical and executed floors are reported side by side.

To accommodate broader clinical study designs, the engine also supports \textbf{paired} designs (one summary per condition per donor with joint within-donor sign flips, yielding $2^{n_{\text{pairs}}}$ allocations) and \textbf{stratified} designs (permuting donor labels within declared batch or slide strata while preserving stratum-specific group counts). For paired inference the statistic is the mean within-donor condition difference after feature-wise scaling by the marginal standard deviation of the pair-centered observations. Its permutation distribution uses joint sign flips across features. Because standard FDR guarantees require valid individual $p$-values and suitable dependence structure, public-data BH call fractions are interpreted as panel sensitivity summaries rather than realized false-discovery proportions.

\subsection{Dimension 2: Acquisition-design and stage collinearity}
To assess whether clinical groups are confounded with categorical acquisition factors (\texttt{slide\_id} or \texttt{batch}), patient weights are normalized to sum to one per donor before computing normalized mutual information (NMI). Binary independent contingency tables receive a Fisher exact test, while multi-level $R \times 2$ tables receive a fixed-margin Monte Carlo permutation test. When metadata is partially missing, the engine logs the missing fraction and tests whether missingness itself associates with group membership.

On physical stage coordinates (or local section bounding-box centers, explicitly labeled when a shared acquisition frame is unavailable), linear feasibility programming detects complete linear separation. Alongside descriptive in-sample AUC, the engine computes patient-level centroid separation and leave-one-patient-out cross-validated AUC (\texttt{lopo\_cv\_auc}) to mitigate finite-sample optimism, together with a sparse least-squares estimability check across fully recorded acquisition main effects. Constant patient centroids give an undefined cross-validated AUC. Disjoint stage support signals extrapolation. Exact aliasing requires the group indicator to lie in the span of the measured fixed effects.

\subsection{Dimension 3: Pseudoreplication sensitivity}
For every registered feature, the engine contrasts a pooled observation-level two-sided tie-corrected Mann--Whitney test against a donor-level moderated permutation test. In the donor-level test, section means are averaged equally within each biological individual; each feature is standardized by its label-blind marginal standard deviation, and feature variances are moderated with a scaled inverse-chi-squared prior estimated via log-moment matching \cite{limma}. Crucially, the variance prior is refit under every label allocation in vectorized batches, ensuring symmetric algorithmic treatment across all permutations while clinical exchangeability remains an explicit design assumption. For nonnegative count measurements, the engine additionally computes a library-size normalized patient contrast and group depth ratio to separate sequencing-depth shifts from sampling-unit effects.

\subsection{Dimension 4: Multi-scale graph diffusion and covariate projection}
By default, occupied square control volumes of width $h$ ($\mu\mathrm{m}$) define a no-flux finite-volume graph Laplacian $L = (D - A)/h^2$ that explicitly encodes tissue occupancy and discretization. For irregular single-cell point clouds, an optional clipped Voronoi finite-volume graph mode provides rigid-motion (rotation and translation) invariance. Across $K \ge 1$ user-specified or automatically recommended physical wavelengths (enforcing connected 2-D footprint support and $\lambda \ge \pi h$), a shared Chebyshev polynomial expansion evaluates difference-of-heat-kernels bandpass filters \cite{wavelets}.

Within each band, filtered molecular signals are projected onto an orthonormal basis spanning the filtered covariates; band energies are averaged equally within and across matched patients before computing the arithmetic mean of feature-wise explained energy fractions. To guard against over-parameterized projections, the engine estimates effective spectral dimension and white-noise transfer gain using exact evaluation for small grids ($n \le 144$) and 48 seeded Rademacher probes with standard errors for larger grids. For MALDI-MSI, an optional leave-one-ion-out TIC mode subtracts each target mass window from its own TIC covariate before projection.

\subsection{Dimension 5: High-frequency short-lag noise sensitivity}
Fine-band residual energies are adjusted using a nonnegative matched-path two-lag semivariogram intercept---the path average of $(x_a - 2x_m + x_b)^2 / 6$ along aligned three-volume paths---multiplied by the projected filter's white-noise gain. Signed adjusted energies are retained in donor-level comparisons to avoid truncation bias under unequal precision. To examine whether variance stabilization improves calibration on count data, the engine additionally implements an optional Anscombe transformation mode (\texttt{nugget\_mode="vst"}) prior to fine-band nugget estimation. The section field \texttt{two\_step\_edges} counts matched, equally spaced three-point paths on fine-band support. Irregular Voronoi sites may have no such paths, leaving the nugget unassessed. Convex-hull clipping can bridge unsampled tissue gaps.

\subsection{Cross-language R benchmarking, claim checking, and informatics tools}
Our isolated subprocess harness invokes original R implementations of \texttt{edgeR}, \texttt{DESeq2}, \texttt{DESpace2}, \texttt{nnSVG}, \texttt{SPARK-X}, and \texttt{Cardinal} whenever a cohort satisfies each package's contract \cite{edger,deseq2,despace2,nnsvg,sparkx,cardinal}. Full-library depth is recorded whenever available, and every reported BH call count is independently reconstructed from raw $p$-values over the registered feature family. Complementing the core engine, \texttt{claim-check} extracts bilingual unit counts, group-count pairs, and negation-aware statistical strategy mentions from methods text or PDFs; \texttt{spatial-audit-mcp} exposes seven local stdio tools for AI-assisted workflows; and \texttt{docs/index.html} delivers a zero-backend interactive observatory with searchable per-feature diagnostics, SVG chart export, local JSON upload, and an interactive English/Chinese interface.

\section{Results}
\subsection{A curated multi-platform human spatial archive}
Our evidence catalog documents archive accessions, clinical contrasts, unit-mapping caveats, and---for all 23 studies with a linked article---the primary publication and scoped methods passage (the supplementary specimen record GSE346380 explicitly lacks a linked publication or named inferential test). The primary human benchmark comprises \ArchiveCount{} distinct accessions across \MeasuredPlatforms{} platform families (11 Visium, five Xenium, and one each for Stereo-seq, CosMx SMI, MERFISH and MALDI-MSI), yielding \PublicRuns{} measured public-data audits (including explicitly labeled panel and donor subsets); cross-study participant overlap cannot be excluded from public metadata alone. Four supplementary entries preserve provenance for cohorts outside the primary independent-donor benchmark: the axolotl injury time series STDS0000056 \cite{axolotl}, a lung CosMx archive with unresolved donor independence, the psoriasis Visium archive GSE225475 lacking a verified biopsy-to-patient crosswalk \cite{geo225475}, and the longitudinal basal cell carcinoma series GSE291246 requiring paired inference \cite{geo291246}. A targeted screening of three additional human CosMx series (GSE348495, GSE312757, GSE311636) found specimen or assay labels without a verifiable independent two-group donor crosswalk in the inspected records and retained them as documented screening records rather than unverified audits (Supplementary Note~7).

The MERFISH series inferred one donor per GEO sample from the published cohort design, distinct demographic tuples and control labels. For fibrosis, the adapter selected the lexicographically first archived core per mapped donor. Table~\ref{tab:coverage} summarizes evaluability by primary accession. Table~\ref{tab:cohorts} retains the separate panel and donor subsets.

\begin{table}[htbp]\centering\small
\caption{Dimension coverage in frozen public reports and additional Wave 1 software capabilities. Each primary archive is counted once.}\label{tab:coverage}
\begin{tabularx}{\linewidth}{@{}p{.17\linewidth}X X p{.18\linewidth}@{}}\toprule
Dimension & Required inputs & Frozen reports (20 archives) & Wave 1 options\\\midrule
Resolution & Donor IDs, arm labels and panel size & 20/20 with permutation floors & Paired and stratified designs\\
Design & Group plus slide, batch or stage metadata. Local centers provide a proxy & 20/20 screened, 0/20 with observed shared stage coordinates & Patient-centroid LOPO AUC\\
Pseudoreplication & Observation values, section-to-donor mapping and group & 20/20 with pooled and donor calls & Count-depth source\\
Composition & Calibrated sites, wavelength support and measured covariates & 11/20 with at least one estimable band & Voronoi graph\\
Nugget & Fine-band support, matched three-point paths and donor features & 9/20 with at least one evaluable feature & Anscombe VST\\
\bottomrule\end{tabularx}
\end{table}
The Wave 1 options (\texttt{paired}, \texttt{stratified}, \texttt{lopo\_cv\_auc}, \texttt{depth\_source}, \texttt{voronoi} and \texttt{vst}) describe available software. Coverage counts use the frozen public JSON reports with their recorded designs and geometry.

\begin{table}[htbp]\centering\small
\caption{Public audit reports. Units are donors. Floor is the exhaustive two-sided permutation floor. Spot and unit columns give pooled-observation and donor-level BH call percentages.}\label{tab:cohorts}
\begin{tabular}{lrrrrrr}\toprule
Public input & Units & Positions & Features & Floor & Spot \% & Unit \%\\\midrule
GSE189184~\cite{geo189184} & 9 & 13,382 & 24 & 0.0159 & 100.0 & 0.0 \\
GSE190732~\cite{geo190732} & 6 & 16,716 & 32 & 0.1 & 100.0 & 0.0 \\
GSE202011~\cite{geo202011} & 9 & 7,392 & 32 & 0.0238 & 100.0 & 0.0 \\
GSE206391~\cite{geo206391} & 8 & 12,619 & 32 & 0.0357 & 90.6 & 0.0 \\
GSE220442-full-panel~\cite{chenAD} & 6 & 25,293 & 2512 & 0.1 & 93.3 & 0.0 \\
GSE220442~\cite{chenAD} & 6 & 25,293 & 32 & 0.1 & 87.5 & 0.0 \\
GSE250346~\cite{geo250346} & 35 & 1,279,572 & 32 & 2.83e-08 & 90.6 & 28.1 \\
GSE267226~\cite{geo267226} & 5 & 21,301 & 29 & 0.2 & 96.6 & 0.0 \\
GSE282123~\cite{geo282123} & 17 & 401,318 & 32 & 8.23e-05 & 93.8 & 9.4 \\
GSE284005~\cite{geo284005} & 17 & 401,794 & 32 & 0.00294 & 93.8 & 0.0 \\
GSE294965~\cite{geo294965} & 12 & 846,950 & 32 & 0.00216 & 100.0 & 56.2 \\
GSE306288~\cite{valveCAVD} & 6 & 12,458 & 25 & 0.1 & 96.0 & 0.0 \\
GSE318039~\cite{lungCOPD} & 8 & 19,788 & 28 & 0.0357 & 85.7 & 0.0 \\
GSE327879~\cite{geo327879} & 6 & 76,973 & 32 & 0.1 & 96.9 & 0.0 \\
GSE336890~\cite{geo336890} & 13 & 227,587 & 17 & 0.00117 & 100.0 & 0.0 \\
GSE348499-subset~\cite{yanTNBC} & 6 & 319,958 & 32 & 0.1 & 96.9 & 0.0 \\
GSE348499~\cite{yanTNBC} & 44 & 3,251,882 & 32 & 1.42e-12 & 84.4 & 0.0 \\
MTBLS12782~\cite{msiPDAC} & 10 & 96,401 & 32 & 0.00794 & 100.0 & 0.0 \\
S-BSST1410~\cite{lungHuman} & 8 & 101,907 & 32 & 0.0286 & 100.0 & 68.8 \\
STT0000036~\cite{crcStereo} & 6 & 78,289 & 32 & 0.133 & 84.4 & 0.0 \\
zenodo-10012934~\cite{lungIPF} & 7 & 57,787 & 25 & 0.0571 & 100.0 & 0.0 \\
zenodo-6578047-FZ~\cite{heartMI} & 7 & 19,835 & 27 & 0.0571 & 100.0 & 0.0 \\
\bottomrule

\end{tabular}
\end{table}

Across all \MeasuredArchives{} measured archives, Figures~\ref{fig:exact-floor}--\ref{fig:nugget} summarize exact permutation resolution, pooled-versus-donor discovery shifts, multiscale covariate shares, and short-distance noise adjustments. Small clinical donor counts frequently placed the exhaustive permutation floor above or near $\alpha = 0.05$, while acquisition metadata directly constrained identifiability: in the CosMx severe COVID-19 lung contrast \cite{geo327879}, case and reference specimens occupy distinct slides, so introducing slide fixed effects absorbs the clinical group contrast entirely.

\auditfigure{exact-floor}{Patient-level permutation resolution across public contrasts.}{Each point marks the theoretical two-sided permutation floor under exhaustive enumeration for the observed biological group sizes. When the combinatorial space exceeds the enumeration budget, the executed Monte Carlo floor is recorded separately in each JSON report. The dashed line indicates the nominal $\alpha=0.05$ threshold.}
\auditfigure{replication}{Sensitivity of discovery rates to the unit of analysis.}{Benjamini--Hochberg call fractions are computed over each cohort's declared feature panel under pooled-observation rank testing versus patient-level moderated permutation testing. Exact enumeration and finite Monte Carlo runs are distinguished in the machine-readable reports.}
\auditfigure{composition}{Covariate share of spatial signal energy across physical wavelengths.}{Each point summarizes patients and sections with valid physical footprint support and non-saturated covariate projections at the indicated micrometre wavelength; denominators can vary across bands and are recorded per cohort. Cortical anatomical labels quantify regional composition, and assay-specific covariates carry their own biological or technical interpretation.}
\auditfigure{nugget}{Fine-scale contrast shifts following short-distance noise adjustment.}{Each point represents one feature normalized by its larger pre-adjustment group energy, with groups ordered lexically as recorded in the report. Departure from the diagonal indicates sensitivity of the fine-band contrast to estimated short-lag white noise.}

\FloatBarrier
\subsection{Clinical contrasts, conditional randomization and validation}
Figure~\ref{fig:case-studies} illustrates three cases from cortex, breast cancer and pancreatic cancer. The fibrosis comparison examines agreement across methods. Table~\ref{tab:null} reports both the number of feature-splits called and the number of splits with any calls. Supplementary Note~7 gives the sampling and validation procedures.

In 27 simulation settings with no batch alignment, each mode used the same count realizations (192 cohorts per effect state and setting). Across settings, the largest null rejection fractions were 10.9\% on the raw scale and 38.5\% after Anscombe transformation. The latter was 74/192 cohorts (pointwise 95\% Wilson interval 31.9\% to 45.6\%) under Poisson sampling with seven donors per arm and a twofold case baseline. At the same donor count and mean shift with dispersion 0.35, alternative detection was 56.2\% and 24.5\%, respectively. A fixed additive pattern in counts changes amplitude after a square-root transformation as baseline abundance changes. These rates evaluate the original latent-amplitude null.

GSE348499 supplied 50 random donor splits within each pretreatment breast cancer response group. Across splits and features, the BH call fractions for pooled cells, donor means, raw-scale corrected energy and VST corrected energy, respectively, were 93.31\%, 2.69\%, 0.12\%, 0.00\% within RD (19 donors) and 95.06\%, 1.69\%, 1.62\%, 1.62\% within pCR (25 donors). Splits reuse donor material. The fractions describe random-label sensitivity conditional on these donors.

In GSE220442, expanding the same 6 cortical donors from 32 to 2,512 archived features preserved all five screening categories. Pooled BH calls rose from 87.5\% to 93.3\%. Both panels had zero donor-level calls. Cortical-layer covariates accounted for 35.9\% and 33.9\% of energy at 1,400 micrometres. Neither panel supported the 250-micrometre nugget endpoint.

For cortical SDF4, the pooled adjusted p was $3\times 10^{-56}$ and the donor-adjusted p was 1. Three donors per arm limit the smallest two-sided donor p to 0.10. In the breast cancer panel, ABAT yielded pooled and donor-adjusted p values of $1.6\times 10^{-30}$ and 0.992. For ABCB4, the pCR-minus-RD energy difference changed from $-3.68\times 10^{-5}$ to $-7.98\times 10^{-6}$ after short-lag subtraction. The ABCB4 contrast compares pretreatment spatial variation between eventual response groups. Without measured cell-type covariates in these sections, the contribution of cellular composition to this difference remains unresolved.

In pancreatic cancer, the without-minus-with-diabetes energy difference for m/z 178.00810 changed from 0.743 to 0.010. The archive identifies this feature by mass window. Across the panel, TIC and raster coordinates together accounted for 45.09\% of energy at 1,400 micrometres, versus 44.88\% after subtracting each target window from its own TIC covariate. The resulting share still includes spatial gradients shared by tissue and acquisition.

The 35-donor lung fibrosis Xenium contrast GSE250346 \cite{geo250346} yielded 9/32 donor-level BH calls, 13 edgeR calls and 16 DESeq2 calls. All three procedures called 9 shared features (ACKR1, AGER, APLNR, ATF6, AXIN2, BCL2L11, BMP4, BMPR2, CA4). AGER and CA4 link the panel to alveolar and vascular compartments, respectively. Whole-tissue summaries combine changes in compartment abundance and expression within cells. All three analyses share donors and assay data.

Across 4 independently launched CPU audits with 10 donors and 64 sparse features, wall time increased from 2.7 to 34.5 seconds and peak RSS from 179.2 to 283.8 MiB between 10,000 and 1,000,000 observations. Each size has one host-specific measurement.


\auditfigure{case-studies}{Feature contrasts from cortex, breast cancer, and pancreatic cancer.}{Left: adjusted probabilities for cortical \textit{SDF4} (dashed line at 0.05). Center and right: group differences in fine-band energy before and after nugget subtraction. Scales are assay-specific; an archived mass window does not establish chemical identity.}

The factorial simulation comprised \PhaseCells{} design states and \SimCohorts{} cohorts. Among settings with zero batch alignment and at least five donors per arm, the maximum adjusted null rejection was \MaxUnalignedNull{}\%. The maximum on the prespecified slice in Figure~\ref{fig:phase-diagram} was \MaxNullSlice{}\%. Figure~\ref{fig:vst-scaling} shows the matched variance-stabilization experiment and resource measurements. Table~\ref{tab:scaling} gives the four workload sizes.

\begin{center}\small\begin{tabular}{llrrr}\toprule
Arm & Endpoint & Nominal (\%) & BH calls/tests & Splits with calls\\\midrule
RD & Pooled cells & 93.44 & 1493/1600 & 50/50\\
RD & Donor mean & 7.94 & 43/1600 & 4/50\\
RD & Raw-scale energy & 6.25 & 2/1600 & 2/50\\
RD & VST energy & 6.62 & 0/1600 & 0/50\\
pCR & Pooled cells & 95.06 & 1521/1600 & 50/50\\
pCR & Donor mean & 3.56 & 27/1600 & 1/50\\
pCR & Raw-scale energy & 3.94 & 26/1600 & 2/50\\
pCR & VST energy & 3.62 & 26/1600 & 2/50\\
\bottomrule\end{tabular}\end{center}
\begin{center}\small\begin{tabular}{rrr}\toprule
Observations & Audit seconds & Peak RSS (MiB)\\\midrule
10,000 & 2.69 & 179.2\\
50,000 & 8.94 & 187.8\\
250,000 & 16.65 & 203.0\\
1,000,000 & 34.52 & 283.8\\
\bottomrule\end{tabular}\end{center}


\auditfigure{phase-diagram}{Conditional operating characteristics across the synthetic design grid.}{Panels display adjusted null rejection and alternative-hypothesis detection across the prespecified parameter slice recorded in the simulation JSON. The maximum adjusted null rejection on this slice is \MaxNullSlice{}\%. Surface interpolation between evaluated grid cells is for visual clarity.}

The benchmark comprised \SOTARuns{} native method--cohort runs with independently reconstructed BH calls (Table~\ref{tab:comparators}). The remaining 43 method--cohort pairs were recorded as inapplicable. Figure~\ref{fig:benchmark} shows measured time and memory by method.

\begin{table}[htbp]\centering\small
\caption{Native R method coverage across 20 registered primary contrasts. Multiple methods can run on the same cohort.}\label{tab:comparators}
\begin{tabular}{@{}llr@{}}\toprule
Method & Target & Completed runs\\\midrule
edgeR & Donor pseudobulk differential expression & 18\\
DESeq2 & Donor pseudobulk differential expression & 18\\
DESpace2 & Group-by-domain interaction & 3\\
nnSVG & Within-section spatial variability & 19\\
SPARK-X & Within-section spatial variability & 18\\
Cardinal & Donor-averaged pixel intensities & 1\\
\bottomrule\end{tabular}
\end{table}

\auditfigure{vst-scaling}{Calibration under an abundance shift and observed computational scaling.}{Top: exhaustive donor tests with seven donors per arm and a twofold baseline shift; bars give pointwise 95\% Wilson intervals across independent cohorts. Bottom: one fresh CPU audit per observation count at fixed panel size and spatial extent.}
\auditfigure{benchmark}{Observed runtime and memory across original-package R comparators.}{Each point corresponds to one verified method-by-cohort subprocess with measurable runtime and peak process-tree resident memory, colored by method. Measurements include R startup, matrix ingestion, and model fitting.}

\section{Discussion}
\subsection{Practical guidance for clinical spatial omics workflows}
From a biomedical informatics perspective, the five audit dimensions translate directly into actionable pre-publication checks:
\begin{enumerate}
\item \textbf{Power red/yellow flag:} When the exact permutation floor exceeds $\alpha$ or the first feasible BH rank is large, report donor-level effect sizes descriptively or expand biological replication rather than substituting cell-level $p$-values.
\item \textbf{Acquisition design:} Exact fixed-effect aliasing prevents separate estimation of clinical and acquisition effects. Collect observations under a crossed design. Geometrically disjoint stage regions signal reliance on extrapolation even when a linear group coefficient is estimable. Inspect spatial support and acquire overlapping positions where feasible.
\item \textbf{Pseudoreplication red flag:} When pooled spot discovery rates greatly exceed donor-level rates, inspect library-size ratios and donor-level trajectories before prioritizing candidate biomarkers.
\item \textbf{Composition red/yellow flag:} High within-band covariate shares indicate that spatial energy aligns strongly with anatomical regions, cell-type proportions, or TIC; interpret downstream spatial claims conditional on tissue architecture and consider leave-one-ion-out TIC for MSI.
\item \textbf{Nugget flag:} Treat flagged fine-scale features as abundance-noise-sensitive candidates requiring higher-resolution validation rather than definitive differential spatial markers.
\end{enumerate}

\subsection{Limitations and scope}
Several methodological boundaries warrant explicit emphasis. First, public-cohort inclusion depended on verifiable donor crosswalks, calibrated micrometre coordinates, and an independent two-group contrast; this convenience collection cannot estimate the literature-wide prevalence of design issues, and MALDI-MSI, CosMx, MERFISH and Stereo-seq each contribute a single primary archive. Second, pooled versus patient-level BH call differences on public data reflect concurrent shifts in sampling unit and estimand (pooled rank distributions versus donor means) and do not directly estimate realized false-discovery proportions. Third, as demonstrated by both our factorial grid and our Anscombe VST evaluation, neither raw nor variance-stabilized short-lag subtraction guarantees nominal Type~I error under strong abundance-dependent heteroscedasticity; signed nugget adjustment remains a diagnostic screen rather than a calibrated inferential test. Fourth, observational clinical contrasts still require an external exchangeability argument beyond any software check, and fixed targeted panels cannot rule out unmeasured biology. Evaluation with clinicians, pathologists and study statisticians is future work, including whether the audit changes marker selection, requests for new samples or interpretation of a negative result.

\section{Conclusion}
\textit{Spatial Audit} unifies five essential statistical and experimental-design diagnostics into an open-source, reproducible biomedical informatics resource spanning six spatial omics platforms. By pairing exact and stratified patient-level randomization, acquisition-collinearity verification, physical-scale covariate orthogonalization, real-data empirical-null benchmarking, and a zero-backend interactive observatory, the framework equips researchers and reviewers to evaluate the statistical foundation of multi-patient spatial studies transparently and efficiently.

\section*{Data and Code Availability}
The MIT-licensed \textit{Spatial Audit} package, archive preparation scripts, R comparator harness, simulation suite, local MCP server, offline interactive dashboard, and manuscript build pipeline are publicly available at \url{https://github.com/XG221B/spatial-audit} (live observatory: \url{https://xg221b.github.io/spatial-audit/}). All measured public audit reports link standardized inputs to upstream repositories via SHA-256 receipts, and underlying clinical datasets remain governed by their respective archive terms.

\section*{Declarations}
\noindent\textbf{Ethics:} This secondary analysis used de-identified human data made available through public archives. The author had no contact with participants, collected no new specimens and did not attempt re-identification. Original studies describe their consent and ethics procedures.

\noindent\textbf{Competing interests:} The author declares no competing interests. \textbf{Funding:} This work received no external funding. \textbf{AI usage disclosure:} Large language models and agentic coding assistants were used to assist with software engineering, test generation, documentation, and manuscript editing; all statistical logic, numerical outputs, archive receipts, and release-gate checks were verified through automated testing, cross-language R reconstruction, and author review.

\nocite{*}
\bibliographystyle{unsrt}
\bibliography{references}
\end{document}
